\section{Literature: Predicting NFL Games and the Efficiency of the Lines}\label{sec:literature}

The first question raised after the 30 September report was whether anyone has tried to predict NFL scores before, what data they used, and how they thought about the problem. The literature falls into five groups.

\paragraph{Rating models from scores alone.} Harville~\cite{harville1980} fitted a mixed linear model to score differences from 1,320 games and obtained an average error of 10.68 points against 10.49 for the bookmaker. Glickman and Stern~\cite{glickmanstern1998} let each team's strength drift from week to week and season to season in a Bayesian state-space model and matched the Las Vegas line on a 1993 test set. Their handbook chapter~\cite{glickmanstern2017} surveys these methods and recommends entering each covariate as the difference and the average of the two teams' values, which we adopted. FiveThirtyEight's Elo system~\cite{fivethirtyeightelo} is the practical version: one strength number per team updated from margin of victory, with adjustments for home field, rest and the starting quarterback. All of these use only results and the schedule, and all land a few tenths of a point behind the line.

\paragraph{Exact scores and distributions.} Baker and McHale~\cite{bakermchale2013} is the only paper that forecasts both teams' exact NFL scores, with a birth-process model and pre-game covariates; it too performs about as well as the bookmaker. Stern~\cite{stern1991} showed that margin minus spread is approximately normal with standard deviation near 14, which gives the standard conversion from a spread to a win probability; Stern~\cite{stern1994} and Polson and Stern~\cite{polsonstern2015} extend this to the progress of a game and to a market-implied volatility. In soccer, Dixon and Coles~\cite{dixoncoles1997} model scores as correlated Poisson counts and Divos et al.~\cite{divos2018} price in-play bets by risk-neutral methods; the lesson for us is that football scoring is discrete in a way Poisson models do not capture.

\paragraph{Play-level information.} Yurko, Ventura and Horowitz~\cite{yurko2019} defined expected points and win probability from play-by-play data and showed how to attribute them to plays and players. Expected points added per play is the efficiency measure the analytics community has settled on; the nflverse project provides it for every play since 1999 and it is the source of most of our team-form features.

\paragraph{Market efficiency.} Levitt~\cite{levitt2004} showed from a bettor-level dataset that bookmakers do not balance their books but set prices as skilled forecasts that exploit bettor biases, which is why the closing line is hard to beat. Boulier and Stekler~\cite{boulierstekler2003} and Song, Boulier and Stekler~\cite{songboulierstekler2007} compared statistical systems, experts and the betting market and found the market most accurate, with models and experts tied. Zuber, Gandar and Bowers~\cite{zuber1985} and Gray and Gray~\cite{graygray1997} found only small, unstable inefficiencies. Lopez, Matthews and Baumer~\cite{lopez2018} go the other way and use moneylines as the best available measure of team strength. Shin~\cite{shin1993} and Wolfers and Zitzewitz~\cite{wolferszitzewitz2006} show how to read probabilities from odds and from prediction-market prices.

\paragraph{Machine-learning studies.} A systematic review by Galekwa et al.~\cite{mlreview2024} and a number of student projects report accuracies of 70 to 95 percent at picking winners. Where the methods are described, these come from same-game statistics (leakage) or from very small test sets. The closing line picks the winner about 66 percent of the time; nothing credible is far above that. We took two lessons: build an explicit leakage test, and do not expect a flexible learner to add much over a well-specified rating model with 7,000 games.

\paragraph{Data in the literature.} The serious score models use only scores, dates and home field; richer inputs appear in three families: play-level efficiency, quarterback identity and betting lines. Bet-volume data (Levitt; Paul and Weinbach) is proprietary; the Kalshi trade tape is the free equivalent. No paper reviewed uses injuries, snap counts or depth charts in a score model, although Glickman and Stern~\cite{glickmanstern2017} name player-level factors as the obvious gap. Everything the literature uses is available free from nflverse except opening lines and bet splits.

\paragraph{Implication for our design.} Point accuracy for NFL scores is capped near the closing line for every published method. The sensible goal is a calibrated distribution over both scores, consistent across every derived quantity, rather than a lower point error. The suggestions received from Professors Amine and Shasha fit this picture: a multi-output random forest can use every feature without selection, a linear baseline shows whether the forest adds anything, and a probabilistic extension predicting two means, two standard deviations and a correlation is the natural destination for a two-score model.
